\subsection{Suffix Array + Algoritmo de Kasai (LCP)}

\graybold{Preguntas guía}

\begin{itemize}
\item ¿Necesito comparar eficientemente muchos sufijos o subcadenas de un mismo texto (la subcadena repetida más larga, subcadena común más larga entre dos textos, número de subcadenas distintas, etc.)?
\item ¿El problema se puede reformular como "encontrar el máximo prefijo común entre dos sufijos del texto"?
\item ¿La longitud del texto vuelve inviable comparar todas las subcadenas por fuerza bruta, pero admite una solución \bigO{n \log^2 n} o mejor?
\item ¿Las ocurrencias repetidas pueden solaparse entre sí dentro del texto original?
\end{itemize}

\graybold{¿Para qué sirve?} El \emph{suffix array} es el arreglo de todas las posiciones iniciales de los sufijos de un texto, ordenadas lexicográficamente según el sufijo que representan. Junto con el \emph{arreglo LCP} (longest common prefix entre sufijos consecutivos en ese orden, obtenido con el algoritmo de Kasai), forma la base de una familia enorme de algoritmos de cadenas: subcadena repetida más larga, subcadena común más larga entre dos textos, conteo de subcadenas distintas, búsqueda de patrones, entre otros. La idea central es que cualquier propiedad sobre "todas las subcadenas de un texto" se puede reducir a una propiedad sobre "prefijos de sufijos ordenados", lo cual es mucho más manejable.

\graybold{Complejidad:} \bigO{n \log^2 n} para construir el suffix array por duplicación con ordenamiento estándar (se puede bajar a \bigO{n \log n} con radix sort, innecesario para \(n \le 5000\)), más \bigO{n} para el arreglo LCP con el algoritmo de Kasai.

\graybold{Fundamento teórico:} La construcción por duplicación mantiene, en la iteración $k$, un rango $\mathrm{rank}_k[i]$ que ordena correctamente todas las subcadenas de longitud $2^k$ que empiezan en $i$. Para pasar de longitud $k$ a $2k$, basta comparar el par $(\mathrm{rank}_k[i], \mathrm{rank}_k[i+k])$: el primer componente ordena la primera mitad de longitud $k$ y el segundo la segunda mitad, así que ordenar por ese par ordena correctamente la subcadena completa de longitud $2k$ (invariante que se preserva inductivamente, partiendo de $\mathrm{rank}_0$ = valor ASCII de cada carácter). Tras $\lceil \log_2 n \rceil$ iteraciones, $2^k \ge n$ y el rango final es exactamente el orden de los sufijos completos.

El algoritmo de Kasai calcula $\mathrm{lcp}[i]$ (LCP entre el sufijo en la posición $i$ del suffix array y el anterior) recorriendo el texto en su orden original, no en el orden del suffix array. Su observación clave es que si el sufijo que empieza en la posición $p$ tiene LCP $h$ con su vecino en el suffix array, entonces el sufijo que empieza en $p+1$ tiene LCP de al menos $h-1$ con \emph{su} vecino correspondiente (quitar el primer carácter de dos cadenas que compartían $h$ caracteres deja $h-1$ compartidos como cota inferior). Esto permite nunca decrementar $h$ en más de 1 entre iteraciones consecutivas, acotando el trabajo total en \bigO{n} de forma amortizada, en vez de recalcular cada LCP desde cero.

La respuesta al problema (subcadena repetida más larga, con solapamiento permitido) es exactamente $\max(\mathrm{lcp})$: dos sufijos adyacentes en el suffix array con $\mathrm{lcp} = L$ comparten un prefijo común de longitud $L$, que por definición es una subcadena que aparece empezando en dos posiciones distintas del texto original — sin restricción alguna sobre si esas dos ocurrencias se solapan o no.

\graybold{Ejercicio aplicado}

Dado un texto $S$ de hasta 5000 caracteres minúsculos, hallar la longitud de la subcadena más larga que aparece al menos dos veces en $S$ (las dos ocurrencias pueden solaparse).

\graybold{I/O:} Se usó lectura línea a línea (\texttt{sys.stdin.buffer.readline}) porque cada caso es simplemente una línea con la cadena completa, sin necesidad de tokenizar ni detectar separadores especiales. Se midió el rendimiento con un caso de estrés cercano al peor caso ($n=5000$ por cadena): con $T=30$ cadenas de un solo carácter repetido (peor caso para el número de iteraciones de duplicación) tardó 0.84s, y con $T=200$ cadenas de $n=5000$ con alfabetos mixtos (un carácter, dos caracteres, alfabeto completo) tardó 3.72s (~18ms por caso). Dado que este problema típicamente trae pocos casos de prueba, el rendimiento es holgado; si $T$ fuera mucho mayor (miles de casos), convendría reemplazar el ordenamiento por duplicación por radix sort para bajar a \bigO{n \log n}.

\begin{lstlisting}[language=Python]
import sys

def longest_repeated_substring(s):
    n = len(s)
    if n <= 1:
        return 0

    # --- Construccion del suffix array por duplicacion (O(n log^2 n)) ---
    sa = list(range(n))
    rank = [ord(c) for c in s]
    tmp = [0] * n
    k = 1
    while True:
        # La clave de comparacion combina el rango de la primera mitad
        # (longitud k) y el de la segunda mitad, para ordenar subcadenas
        # de longitud 2k a partir del orden ya conocido de longitud k
        def key(i):
            return (rank[i], rank[i + k] if i + k < n else -1)
        sa.sort(key=key)
        tmp[sa[0]] = 0
        for i in range(1, n):
            # Nuevo rango: mismo valor si empatan, +1 si cambia el orden
            tmp[sa[i]] = tmp[sa[i - 1]] + (1 if key(sa[i - 1]) < key(sa[i]) else 0)
        rank = tmp[:]
        if rank[sa[-1]] == n - 1:
            break  # todos los sufijos ya tienen rango distinto: orden total logrado
        k <<= 1

    # --- Algoritmo de Kasai para el arreglo LCP en O(n) amortizado ---
    lcp = [0] * n
    h = 0  # LCP acumulado, nunca baja en mas de 1 entre posiciones consecutivas
    for i in range(n):
        if rank[i] > 0:
            j = sa[rank[i] - 1]  # sufijo vecino en el suffix array
            while i + h < n and j + h < n and s[i + h] == s[j + h]:
                h += 1
            lcp[rank[i]] = h
            if h > 0:
                h -= 1
        else:
            h = 0

    return max(lcp) if lcp else 0


def main():
    data = sys.stdin.buffer
    t = int(data.readline())
    salidas = []
    for _ in range(t):
        s = data.readline().decode().strip()
        salidas.append(str(longest_repeated_substring(s)))
    sys.stdout.write('\n'.join(salidas) + '\n')


if __name__ == "__main__":
    main()
\end{lstlisting}
